\subsection{Explicit structured behavior-tree teaching}

Structured teaching complements, rather than replaces, ordinary action learning.
Ordinary instruction teaching still builds a linear ASL action. It does not infer
arbitrary control structure. An instruction script must explicitly select a
parent node, control kind, stable node identifier, and insertion index.
The structured learner is owned by each agent's existing ActionLearning instance.
It does not intercept ordinary execution goals or add another language interpreter.

The root node has identifier \texttt{root}. Native controls, decorators, and
built-in leaves can be taught with the generic node API:
sequence/fallback variants (including async and reactive variants), parallel
controls, conditionals, \texttt{TryCatch}, \texttt{Switch2}--\texttt{Switch6},
all native decorators (including typed loops and update watchers), and
\texttt{AlwaysSuccess}, \texttt{AlwaysFailure}, \texttt{Script},
\texttt{ScriptCondition}, \texttt{SetBlackboard}, \texttt{UnsetBlackboard},
\texttt{Sleep}, and \texttt{WasEntryUpdated}.
Unknown kinds are rejected, not flattened into a sequence.
Existing ASL actions, including actions produced by ordinary action learning,
are leaves through the existing explicit BT binding format.
Their native conditions, effects, and Belief interaction are unchanged.
No persistent world-state store is introduced.

\paragraph{Instruction services.}
The implemented GoalManager services delegate to
\texttt{em.getBehaviorTreeLearning(actor)}. With \texttt{Symbol actor},
\texttt{String name}, and all other textual fields of type \texttt{String},
the service signatures and learner delegates are:
\begin{verbatim}
boolean beginBehaviorTreeLearning(actor, name, kind, bindingsFile)
  -> beginFromFile(name, actor, kind, bindingsFile)
boolean beginBehaviorTreeLearning(actor, name, kind, bindingsFile, attributesJson)
  -> beginFromFile(name, actor, kind, bindingsFile, attributesJson)
boolean addBehaviorTreeNode(actor, name, parent, id, kind, attributesJson, int index)
  -> addNode(name, parent, id, kind, attributesJson, index)
boolean addBehaviorTreeDefinition(actor, name, treeId, rootId, kind, attributesJson)
  -> addDefinition(name, treeId, rootId, kind, attributesJson)
boolean addBehaviorTreeSubtree(actor, name, parent, id, treeId, remapsJson, int index)
  -> addSubtree(name, parent, id, treeId, remapsJson, index)
boolean addBehaviorTreeControl(actor, name, parent, id, kind, int index)
  -> addControl(name, parent, id, kind, index)
boolean addBehaviorTreeLeaf(actor, name, parent, id, binding, portsJson, int index)
  -> addLeaf(name, parent, id, binding, portsJson, index)
boolean recordBehaviorTreeStep(actor, name, parent, id, binding,
                              Predicate step, int index)
  -> record(name, parent, id, binding, step, index)
boolean removeBehaviorTreeNode(actor, name, id) -> remove(name, id)
boolean moveBehaviorTreeNode(actor, name, id, parent, int index)
  -> move(name, id, parent, index)
boolean endBehaviorTreeLearning(actor, name) -> finish(name)
boolean cancelBehaviorTreeLearning(actor, name) -> cancel(name)
boolean exportBehaviorTree(actor, name, xmlFile, bindingsFile)
  -> export(name, xmlFile, bindingsFile)
String getLearnedBehaviorTreeXml(actor, name) -> getXml(name)
\end{verbatim}
Mutation services have \texttt{@Action} and \texttt{@TRADEService};
the XML query has only \texttt{@TRADEService}. File operations may throw
\texttt{IOException}; malformed edits throw \texttt{IllegalArgumentException}.
These service signatures are excluded from legacy linear goal recording.
Actor lookup uses the existing ExecutionManager agent/team registry; null or
unknown actors are explicit errors rather than creating unowned teaching state.

\paragraph{Complete teaching example.}
After starting GoalManager with \texttt{-asl bt-examples.asl}, the supplied
binding \texttt{MoveBase} resolves to the existing ASL action
\texttt{exampleMoveBase}. The following Java client teaches and publishes a
one-step sequence, exports both artifacts, then executes the learned action.
The output directory must already exist. These same services can be called
from instruction ASL using the action dispatch prefix.
\begin{verbatim}
Symbol actor = Factory.createSymbol("self:agent");
String bindings =
    "core/src/main/resources/config/edu/tufts/hrilab/action/bt/"
    + "example-bindings.json";
manager.beginBehaviorTreeLearning(actor, "taughtVisit",
    "Sequence", bindings);
manager.addBehaviorTreeLeaf(actor, "taughtVisit", "root",
    "visit", "MoveBase", "{\"goal\":\"inspection_station\"}", 0);
manager.endBehaviorTreeLearning(actor, "taughtVisit");
manager.exportBehaviorTree(actor, "taughtVisit",
    "taughtVisit.xml", "taughtVisit-bindings.json");
long id = manager.submitGoal(
    Factory.createPredicate("taughtVisit", actor));
GoalStatus result = manager.joinOnGoal(id);
// A later session can reload the exported pair under a fresh action name:
manager.loadBehaviorTree("taughtVisit.xml",
    "taughtVisit-bindings.json", "reloadedVisit");
\end{verbatim}
The binding path above is relative to the repository root; when launching
from the \texttt{core} directory, omit its leading \texttt{core/}.
To teach fallback instead of a linear sequence, explicitly add a
\texttt{Fallback} control and its ordered alternatives; no such structure is
inferred from the example's single ordinary action call.

\paragraph{Editing and recording.}
Root parameters use the overloaded begin service. For example, a root
\texttt{Repeat} uses \texttt{\{"num\_cycles":"3"\}} as attributes.
The generic node API accepts the same attributes as the native converter,
including scripts, pre/postcondition hooks, thresholds, queue and update ports,
and explicit leaf bindings. Its attributes and child arities are validated
through the converter at export/publication, not during partial editing.
For example, a \texttt{TryCatch} can be created before its catch branch, and
a decorator can be created before its single child.
The existing add-control/add-leaf/record services remain supported.
Explicit callable definitions use \texttt{addBehaviorTreeDefinition}; their roots
are addressed by \texttt{rootId}. The main definition ID remains the taught action
name, and its root node remains \texttt{root}. Every node ID must be unique across
all definitions in the taught model, so existing parent/edit APIs also address
definition bodies without introducing implicit namespaces.
Definition roots cannot be removed or moved. Ordinary nodes can be moved between
definitions; structural parent/child cycles remain invalid.

\texttt{addBehaviorTreeSubtree} adds a retained \texttt{SubTree} call node with an
explicit target \texttt{treeId} and JSON remappings, for example
\texttt{\{"depth":"\{next\}"\}}. Calls can refer forward, directly recurse, or
mutually recurse. Those call-graph cycles are intentional and are distinct from
invalid structural edit cycles. Target references, remappings, and child arities
are validated through the converter at export/publication. Even unused retained
definitions are validated, rather than allowing invalid instructions to hide.
All definitions and call nodes are exported in the same XML without inlining.
External include-file editing is not inferred; define the required trees
explicitly or bind an existing ASL action as a leaf.

Indices are zero-based. Moving a node uses the destination list after removal;
removing a control removes its entire subtree. Duplicate identifiers, cycles,
leaf parents, bad indices, unknown bindings, missing inputs, and unknown ports
are rejected without applying the edit by the explicit leaf API. Generic-node
attributes and required inputs are retained and rejected at export/publication
if invalid. Incomplete branches are permitted while editing but not at export
or publication.
\texttt{portsJson} is a JSON object mapping registered port names to values;
explicit blackboard mappings such as \texttt{\{"label":"\{target\}"\}} are retained.
They remain invocation-local BT/ASL data, not a second representation of world state.

Recording takes an existing action predicate with the actor as its first
argument, matches it against the selected binding, and preserves ordered concrete
arguments using the existing EventSpec convention. It does not infer input roles,
output mappings, actor changes, consultant declarations, or branching.
Instructions requiring role abstraction or output mappings must use explicit
leaf ports. A separately learned ASL action can encapsulate richer existing
learning behavior without losing its structure or metadata.

\paragraph{Publication and persistence.}
Finishing validates through the native converter, rejects conflicting action
names, installs the resulting ASL actions, and seals the retained model.
Failed publication leaves the draft editable.
Both finished trees and complete drafts can be exported as format-4 XML and
the snapshotted bindings JSON. Both files are required for later
\texttt{loadBehaviorTree(xmlFile, bindingsFile, actionName)}.
XML is generated from retained nested structure, not reconstructed from the
flattened generated action database. Finished models cannot be edited or
cancelled; export them before discarding their owning execution manager.
